\SourceLeaf{370}%
2000 years ago. Such a candid and unreserved admission seems to us more in keeping with the dignity of science than the vain effort to hide a gap that cannot be filled behind an untenable web of sham proofs.

The author of the first work made a similar attempt 15 years ago in a short treatise, \textit{Tentamen novae parallelarum theoriae notione situs fundatae}, in which he sought to base everything on the concept of identity of position. He defines parallel lines as straight lines that have the same position. From this he concludes that any third straight line must intersect them at equal angles, since these angles are simply the measure of how far the position of this third line differs from the positions of the two parallel lines. This form of proof is repeated in the present work, and we could hardly say that it has gained strength from the philosophical considerations woven into it. Every geometer must dispute the assertion on p.~24, Notionem situs e geometria adeo non excludi posse, ut potius notionibus eius fundamentalibus annumeranda sit, \textit{dudum omnes agnovere geometrae}, in the sense in which the author uses the concept of position in his proof. If we start from the author's definition, Situs est modus, quo plura coëxistunt vel iuxta se existunt in spatio, position is a purely relational concept. We can certainly say that two straight lines \(A,B\) have a certain position relative to one another that is the same as the relative position of two other lines \(C,D\). In his proof, however, the author uses the word position as an absolute concept, speaking of identity of position for two straight lines that do not coincide. This meaning is clearly empty and without foundation until we know what to understand by such an identity and how to recognize it. If it is to be recognized by equality of angles with \textit{one} third straight line, then without a prior proof we still do not know whether the same equality holds for the angles with a fourth straight line. If equality of angles with \textit{every} other straight line is to be the criterion, then we still do not know whether the same position is possible without coincidence. Thus \textit{after} the author's proof we stand exactly where we stood \textit{before} it.

Much of the work concerns the claim against \textsc{Kant}
